\documentclass{article}
\begin{document}

\section{Defining Creativity}
Creativity is defined as:
> The ability to produce or use original and unusual ideas.
\begin{itemize}
    \item Creativity is a complex and multifaceted concept. Creativity can be high in some aspect(s) but low in others. For example, a portrait of a person is not original but the representation can be (i.e., Picaso)
    \item To "think outside of the box" requires knowledge of the box. Therefore, creativity requires a certain level of contextual knowledge and understanding. This is also dependent on what the norms are of the current time period.
    \item Creativity can exist without being quality or logical. For example, a child can be creative in their play but it may not be logical or of high quality.
\end{itemize}

\section{Current Frameworks}
\begin{describe}
    \item[N-Gram Originality vs. Training Data] The absolute baseline of AI creativity is checking if the model is just memorizing text. Researchers measure this by calculating the fraction of unseen n-grams (sequences of $n$ words or characters) in the AI's output compared to its training corpus.
    \begin{itemize}
        \item If the AI writes a poem and uses 4-word phrases that never appeared anywhere on the internet or in its training data, its structural originality score spikes.
        \item $\text{Originality Score} = 1 - \frac{\text{Count of } n\text{-grams found in training data}}{\text{Total } n\text{-grams in output}}$
    \end{itemize}
    \item[The Creativity Frontier: The Harmonic Mean] To prevent rewarding nonsense, the state-of-the-art approach combines the unsupervised distance/originality score with a Quality Score (usually calculated via an automated task metric or an LLM-as-a-judge panel).
    \begin{itemize}
        \item The Creativity Score is calculated as the harmonic mean of originality and quality:
        \item $\text{Creativity Score} = 2 \times \frac{\text{Originality} \times \text{Quality}}{\text{Originality} + \text{Quality}}$
        \item This way, if an output is highly original but of low quality (nonsense), it won't score well overall.
    \item[Self-Similarity and "Intra-Model" Diversity] This looks at its internal variance. If you give an AI the same creative prompt 50 times with a high temperature (randomness setting), does it keep repeating variations of the same idea, or can it explore entirely different regions of the embedding space?
    \begin{itemize}
        \item You can measure this by generating multiple outputs from the AI and calculating the Average Pairwise Distance between its own answers:
        \item $\text{Diversity} = \frac{2}{k(k-1)} \sum_{i < j} \text{Distance}(\text{Output}_i, \text{Output}_j)$
    \end{itemize}
    \item[SemDis (Semantic Distance Platform)] An open platform designed specifically to automate creativity tasks like the Alternate Uses Task (AUT) (e.g., "Name 20 creative uses for a brick"). It passes the AI's answers through unsupervised space to see if the AI suggests typical uses ("building a wall") or highly distant, novel uses ("using it as a crude sandpaper").
    \item[RAG-Based Novelty Checkers] For high-level creative outputs (like an AI generating a new scientific hypothesis or a movie plot), researchers use a Retrieval-Augmented Generation (RAG) loop.
    \begin{itemize}
        \item The AI generates an output, then a retrieval system checks if similar ideas exist in the training data or the broader internet (i.e. top 50 documents from vector database). If the retrieval system finds close matches, the output is flagged as less creative.
        \item This can be formalized as:
        \item $\text{Novelty Score} = 1 - \frac{\text{Similarity}(\text{Output}, \text{Retrieved Data})}{\text{Max Similarity}}$
    \end{itemize}
\end{describe}

\end{document}